% 数学ノート集・共通テンプレート v2（積分論2026年度第1部の版面を参考）
% A5判。章内で定義・定理・例・注意等の番号を共通化。
% コンパイル：platex main.tex（3回），dvipdfmx main.dvi
%              dvipdfmx math_note_template_v2.dvi
\documentclass[a5paper,10pt,onecolumn]{jsarticle}
\usepackage[dvipdfmx,top=22mm,bottom=20mm,inner=19mm,outer=18mm,headheight=13pt,headsep=6mm]{geometry}
\usepackage{amsmath,amssymb,amsthm,mathtools}
\usepackage{enumitem,fancyhdr,titlesec}
\usepackage[dvipdfmx,bookmarks=true,colorlinks=false]{hyperref}
\usepackage{pxjahyper}

% ----- 差し替える文書情報 -----
\newcommand{\NoteTitle}{Baireのカテゴリー定理と一様有界性原理}
\newcommand{\NoteSeries}{数学ノート集}

% ----- 文字・組版 -----
% 参考PDFの質素な明朝＋Computer Modernの印象を引き継ぎながら，
% A5判で窮屈にならないよう本文の行間を少し確保する．
\linespread{1.13}
\setlength{\parindent}{1zw}
\setlength{\parskip}{0pt}
\setlength{\abovedisplayskip}{6pt plus 2pt minus 2pt}
\setlength{\belowdisplayskip}{6pt plus 2pt minus 2pt}
\setlength{\abovedisplayshortskip}{3pt plus 2pt minus 1pt}
\setlength{\belowdisplayshortskip}{4pt plus 2pt minus 1pt}
\setlist[enumerate]{topsep=3pt,itemsep=1pt,parsep=0pt,leftmargin=2.4zw}
\titleformat{\section}{\normalfont\normalsize\bfseries}{\thesection}{0.8zw}{}
\titlespacing*{\section}{0pt}{2.0ex plus .4ex}{.9ex}
\titleformat{\subsection}{\normalfont\normalsize\bfseries}{\thesubsection}{0.7zw}{}
\titlespacing*{\subsection}{0pt}{1.5ex plus .3ex}{.6ex}
\raggedbottom

% ----- 柱とページ番号：ページ上部、罫線なし -----
\fancypagestyle{mathnote}{%
  \fancyhf{}%
  \fancyhead[L]{\footnotesize \NoteSeries\ （\NoteTitle）}%
  \fancyhead[R]{\footnotesize \thepage}%
  \renewcommand{\headrulewidth}{0pt}%
  \renewcommand{\footrulewidth}{0pt}%
}
\pagestyle{mathnote}

% ----- 節内通し番号の定義・定理群 -----
\newtheoremstyle{mathnotetheorem}
 {4.5pt}{4.5pt}{\normalfont}{}%
 {\bfseries}{}{0.65zw}%
 {\thmname{#1}\thmnumber{ #2}\thmnote{（#3）}}
\theoremstyle{mathnotetheorem}
\newtheorem{theorem}{定理}[section]
\newtheorem{definition}[theorem]{定義}
\newtheorem{lemma}[theorem]{補題}
\newtheorem{proposition}[theorem]{命題}
\newtheorem{corollary}[theorem]{系}
\newtheorem{example}[theorem]{例}
\newtheorem{remark}[theorem]{注意}
\renewcommand{\proofname}{証明}
\renewcommand{\qedsymbol}{$\square$}
\makeatletter
\renewenvironment{proof}[1][\proofname]{%
  \par\pushQED{\qed}\normalfont\topsep4\p@\@plus2\p@\relax
  \trivlist\item[\hskip\labelsep\bfseries #1]\ignorespaces
}{\popQED\endtrivlist\@endpefalse}
\makeatother

% ----- 使うノートで適宜拡張する記号 -----
\newcommand{\R}{\mathbb{R}}
\newcommand{\norm}[1]{\lVert #1\rVert}
\newcommand{\abs}[1]{\lvert #1\rvert}

\begin{document}
\section*{読了後にできること}
本稿を読むための前提は，Banach空間，一様有界性原理，双対ノルム．
本稿の具体的な到達課題は，完備性の仮定がないと一様有界性が破れる例を構成する．

\section{完備距離空間とカテゴリー}

\begin{definition}
位相空間 $X$ の部分集合 $A$ が\emph{至る所疎}であるとは，
$\operatorname{int}(\overline A)=\varnothing$ であることをいう．
可算個の至る所疎な集合の和集合を\emph{第一類集合}，
第一類でない集合を\emph{第二類集合}という．
可算個の開集合の共通部分を $G_\delta$ 集合という．
\end{definition}

\begin{lemma}\label{lem:nowhere}
$A\subset X$ が至る所疎であることと
$X\setminus\overline A$ が稠密開集合であることは同値である．
また $X$ の閉集合 $F$ が至る所疎であることと
$F$ の内部が空であることは同値である．
\end{lemma}
\begin{proof}
$X\setminus\overline A$ が開であることは明らかである．
それが稠密であることは，その補集合である
$\overline A$ が空でない開集合を含まないこと，
すなわち $\operatorname{int}(\overline A)=\varnothing$
と同値である．閉集合の場合には $\overline F=F$ を用いる．
\end{proof}

\begin{theorem}[Baireのカテゴリー定理]\label{thm:baire}
$X$ を空でない完備距離空間とし，
$(U_n)_{n\ge1}$ を $X$ の稠密開集合の列とする．
このとき
\[
 \bigcap_{n=1}^\infty U_n
\]
は $X$ で稠密である．特に空でない．
\end{theorem}
\begin{proof}
空でない任意の開集合 $V\subset X$ に対して
$V\cap\bigcap_nU_n\ne\varnothing$ を示す．
$U_1$ は稠密であるから $V\cap U_1$ は空でない開集合である．
$x_1\in V\cap U_1$ を取り，
$r_1>0$ を十分小さくして
\[
 \overline{B(x_1,r_1)}\subset V\cap U_1,
 \qquad 0<r_1<2^{-1}
\]
とする．これは距離空間では開集合の各点のまわりに，
その閉包がもとの開集合に含まれる開球を取れるためである．

帰納的に $x_n,r_n$ が得られたとする．
$U_{n+1}$ は稠密開なので
$B(x_n,r_n)\cap U_{n+1}$ は空でない開集合である．
そこで
\[
 \overline{B(x_{n+1},r_{n+1})}
 \subset B(x_n,r_n)\cap U_{n+1},
 \qquad 0<r_{n+1}<2^{-(n+1)}
\]
となる $x_{n+1},r_{n+1}$ を取る．
この構成から，$m>n$ なら $x_m\in B(x_n,r_n)$ である．
従って $d(x_m,x_n)<r_n<2^{-n}$ なので
$(x_n)$ はCauchy列である．完備性により $x_n\to x$ となる．
固定した $n$ に対しすべての $m>n$ で
$x_m\in\overline{B(x_n,r_n)}$ だから，
その閉性により
\[
 x\in\overline{B(x_n,r_n)}\subset U_n.
\]
また $x\in\overline{B(x_1,r_1)}\subset V$ となる．
ゆえに $x\in V\cap\bigcap_nU_n$ である．
\end{proof}

\begin{corollary}\label{cor:baire-category}
空でない完備距離空間は第一類ではない．
より一般に，完備距離空間 $X$ の第一類集合
$A\subset X$ は，空でない開集合を含まない．
\end{corollary}
\begin{proof}
$A\subset\bigcup_{n=1}^\infty A_n$，
$\operatorname{int}\overline{A_n}=\varnothing$ とする．
補集合 $U_n=X\setminus\overline{A_n}$ は稠密開である．
任意の空でない開集合 $V$ について
定理\ref{thm:baire}により
$V\cap\bigcap_nU_n\ne\varnothing$ となる．
ここで $\bigcap_nU_n\subset X\setminus A$ だから，
$V\not\subset A$ である．特に $A=X$ は不可能である．
\end{proof}

\begin{corollary}[閉集合による被覆]\label{cor:closedcover}
$X$ を空でない完備距離空間とする．
閉集合 $F_n\subset X$ が
\[
 X=\bigcup_{n=1}^\infty F_n
\]
を満たすなら，ある $N$ に対して
$\operatorname{int}F_N\ne\varnothing$ である．
\end{corollary}
\begin{proof}
すべての $F_n$ の内部が空なら，
補題\ref{lem:nowhere}により各 $F_n$ は至る所疎である．
このとき $X$ が第一類となり，
系\ref{cor:baire-category}に矛盾する．
\end{proof}

\begin{remark}
定理\ref{thm:baire}の結論は，
完備性から単に共通部分が空でないというだけでなく，
任意の空でない開集合と交わることまで主張している．
また，カテゴリーの第一類という概念は測度零とは異なる．
例えば $\mathbb R$ の有理数集合は可算で第一類であるが，
密度の点では $\mathbb R$ 全体に稠密である．
\end{remark}

\section{有界線形作用素の準備}

以下，$\mathbb K$ は $\mathbb R$ または $\mathbb C$ とし，
$X,Y$ を $\mathbb K$ 上のノルム空間とする．

\begin{definition}
線形作用素 $T:X\to Y$ が\emph{有界}であるとは，
ある $C\ge0$ が存在して
$\|Tx\|_Y\le C\|x\|_X$ がすべての $x$ に対して
成り立つことをいう．
有界線形作用素全体を $\mathcal B(X,Y)$ と書き，
\[
 \|T\|=\sup_{\|x\|_X\le1}\|Tx\|_Y
\]
を作用素ノルムという．
\end{definition}

\begin{lemma}\label{lem:cont-bound}
線形作用素 $T:X\to Y$ について，
次は同値である．
\begin{enumerate}
\item $T$ は有界である．
\item $T$ は $0$ において連続である．
\item $T$ は全ての点で連続である．
\end{enumerate}
\end{lemma}
\begin{proof}
有界なら
$\|Tx-Ty\|\le\|T\|\|x-y\|$ だから全点で連続である．
全点で連続なら特に $0$ で連続である．
逆に $0$ で連続とすると，
ある $\delta>0$ に対し
$\|u\|<\delta$ なら $\|Tu\|<1$ となる．
$x\ne0$ に対して $u=\delta x/(2\|x\|)$ とすれば
$\|Tx\|<2\|x\|/\delta$ となり有界である．
\end{proof}

\begin{proposition}\label{prop:operator-norm}
$\mathcal B(X,Y)$ は作用素ノルムに関してノルム空間であり，
\[
 \|Tx\|_Y\le\|T\|\|x\|_X
\]
がすべての $x\in X$ に対して成立する．
さらに $Z$ をノルム空間，
$S\in\mathcal B(Y,Z)$ とすれば
$ST\in\mathcal B(X,Z)$ かつ
$\|ST\|\le\|S\|\|T\|$ である．
\end{proposition}
\begin{proof}
$x=0$ のとき最初の不等式は明らかである．
$x\ne0$ に対して $x/\|x\|$ を用いれば定義から従う．
作用素ノルムの斉次性および三角不等式は，
それぞれ $\|(aT)x\|=|a|\|Tx\|$ と
$\|(S+T)x\|\le\|Sx\|+\|Tx\|$ から得る．
$\|T\|=0$ ならすべての $x$ に対して $Tx=0$ である．
最後に
\[
 \|STx\|\le\|S\|\|Tx\|
 \le\|S\|\|T\|\|x\|
\]
を用いる．
\end{proof}

\begin{proposition}[作用素空間の完備性]\label{prop:operator-complete}
$Y$ がBanach空間なら $\mathcal B(X,Y)$ は
作用素ノルムに関してBanach空間である．
ここでは $X$ の完備性は仮定しない．
\end{proposition}
\begin{proof}
$(T_n)$ を $\mathcal B(X,Y)$ のCauchy列とする．
任意の $x\in X$ に対して
\[
 \|T_nx-T_mx\|\le\|T_n-T_m\|\,\|x\|
\]
だから $(T_nx)$ は $Y$ のCauchy列である．
$Y$ の完備性により
$Tx=\lim_nT_nx$ と定めることができる．
極限を取れば $T$ は線形である．

$(T_n)$ は作用素ノルムについて有界なので，
$M=\sup_n\|T_n\|<\infty$ と置ける．
従って任意の $x$ に対して
\[
 \|Tx\|=\lim_n\|T_nx\|\le M\|x\|
\]
であり，$T$ は有界である．
$\varepsilon>0$ に対して $n,m\ge N$ なら
$\|T_n-T_m\|<\varepsilon$ とする．
$m\to\infty$ とすれば
$\|(T_n-T)x\|\le\varepsilon\|x\|$ である．
$\|x\|\le1$ で上限を取って
$\|T_n-T\|\le\varepsilon$ を得る．
\end{proof}

\section{一様有界性原理}

以下，作用素族の添字集合 $\Lambda$ は空でないものとする．

\begin{definition}
作用素族 $(T_\lambda)_{\lambda\in\Lambda}
\subset\mathcal B(X,Y)$ が\emph{各点有界}であるとは，
すべての $x\in X$ に対して
\[
 \sup_{\lambda\in\Lambda}\|T_\lambda x\|_Y<\infty
\]
が成り立つことをいう．
一方，$\sup_\lambda\|T_\lambda\|<\infty$
が成り立つとき，作用素ノルムについて
\emph{一様有界}であるという．
\end{definition}

\begin{theorem}[一様有界性原理，Banach--Steinhaus]\label{thm:ubp}
$X$ をBanach空間，$Y$ をノルム空間，
$(T_\lambda)_{\lambda\in\Lambda}\subset\mathcal B(X,Y)$
を各点有界な作用素族とする．
このとき
\[
 \sup_{\lambda\in\Lambda}\|T_\lambda\|<\infty
\]
である．
\end{theorem}
\begin{proof}
各正整数 $n$ に対して
\[
 F_n=\left\{x\in X:
   \sup_{\lambda\in\Lambda}\|T_\lambda x\|\le n\right\}
 =\bigcap_{\lambda\in\Lambda}
    \{x\in X:\|T_\lambda x\|\le n\}
\]
と置く．各 $T_\lambda$ は連続なので
$F_n$ は閉集合である．
各点有界性により
$X=\bigcup_{n=1}^\infty F_n$ となる．
系\ref{cor:closedcover}から，ある $N$ について
$\operatorname{int}F_N\ne\varnothing$ である．
従ってある $x_0\in X$ と $r>0$ に対して
$B(x_0,r)\subset F_N$ となる．
特に $x_0\in F_N$ だから
\[
 \|T_\lambda x_0\|\le N
 \qquad(\lambda\in\Lambda)
\]
である．$\|u\|<r$ なら $x_0+u\in F_N$ なので
\[
 \|T_\lambda u\|
 \le\|T_\lambda(x_0+u)\|+\|T_\lambda x_0\|
 \le2N
\]
となる．$\|x\|\le1$ とし $u=(r/2)x$ を代入すれば，
すべての $\lambda$ について
\[
 \|T_\lambda x\|\le \frac{4N}{r}.
\]
従って $\sup_\lambda\|T_\lambda\|\le4N/r<\infty$ である．
\end{proof}

\begin{remark}
定理\ref{thm:ubp}は値域 $Y$ の完備性を必要としない．
また，作用素の添字集合 $\Lambda$ は可算でなくてもよい．
証明では添字に関する無限個の閉集合の共通部分を取るが，
Baireのカテゴリー定理を使うのは数列 $(F_n)_{n\ge1}$
に対してである．
\end{remark}

\begin{proposition}[第二類集合上の評価]\label{prop:second-category}
$X$ をBanach空間，$Y$ をノルム空間，
$(T_\lambda)\subset\mathcal B(X,Y)$ とする．
各点有界性が $X$ 全体でなく，
ある第二類集合 $E\subset X$ 上でのみ成立すると仮定しても，
作用素ノルムは一様有界である．
\end{proposition}
\begin{proof}
上の証明で定めた閉集合 $F_n$ に対して
$E\subset\bigcup_nF_n$ が成立する．
もしすべての $F_n$ が内部を持たなければ
各 $F_n$ は至る所疎となり，
$E$ は第一類になってしまう．
よってある $F_N$ は空でない内部を持つ．
後は定理\ref{thm:ubp}の証明と同じ評価を行えばよい．
\end{proof}

\section{各点収束する作用素列}

\begin{theorem}[各点極限の連続性]\label{thm:pointwise-limit}
$X$ をBanach空間，$Y$ をノルム空間とする．
$T_n\in\mathcal B(X,Y)$ が各 $x\in X$ において
$T_nx\to Tx\in Y$ を満たすとする．
このとき $T:X\to Y$ は有界線形作用素であり，
\[
 \|T\|\le\sup_n\|T_n\|<\infty
\]
である．
\end{theorem}
\begin{proof}
任意の $x,y\in X$ および $a,b\in\mathbb K$ に対して
\[
 T(ax+by)=\lim_nT_n(ax+by)
 =aTx+bTy
\]
なので $T$ は線形である．
各 $x$ に対して $(T_nx)$ は収束するから有界である．
定理\ref{thm:ubp}により
$M=\sup_n\|T_n\|<\infty$ となる．
そこで
\[
 \|Tx\|=\lim_n\|T_nx\|
 \le M\|x\|
\]
となり，$T$ は有界である．
\end{proof}

\begin{remark}
定理\ref{thm:pointwise-limit}において，
$Y$ がBanach空間である必要はない．
各点での極限 $Tx$ が $Y$ の元として存在することを
明示的に仮定しているからである．
また，各点収束から作用素ノルム収束は従わない．
\end{remark}

\begin{example}[各点収束とノルム収束の相違]\label{ex:coordinate}
$X=\ell^2(\mathbb K)$ とし
$T_n:X\to\mathbb K$ を $T_nx=x_n$ と定める．
このとき $T_nx\to0$ がすべての $x\in X$ に対して
成り立つが，$\|T_n\|=1$ は全ての $n$ で成り立つ．
\end{example}
\begin{proof}
$\sum_{n=1}^\infty|x_n|^2<\infty$ なら
$x_n\to0$ である．また $|x_n|\le\|x\|_2$ だから
$\|T_n\|\le1$ となる．標準基底の第 $n$ 項 $e_n$ に対し
$\|e_n\|_2=1$ かつ $T_ne_n=1$ であるから
$\|T_n\|\ge1$ である．
\end{proof}

\begin{theorem}[別々の連続性から同時連続性へ]\label{thm:bilinear}
$X$ をBanach空間，$Y$ をノルム空間とし，
$b:X\times Y\to\mathbb K$ を双線形写像とする．
各 $y\in Y$ に対して $x\mapsto b(x,y)$ が連続であり，
各 $x\in X$ に対して $y\mapsto b(x,y)$ が連続なら，
ある $C\ge0$ が存在して
\[
 |b(x,y)|\le C\|x\|\,\|y\|
\]
が全ての $x,y$ に対して成立する．
従って $b$ は積ノルム位相に関して同時連続である．
\end{theorem}
\begin{proof}
$\|y\|\le1$ を満たす各 $y$ に対して
$T_y:X\to\mathbb K$ を $T_y(x)=b(x,y)$ と置く．
仮定から $T_y\in X^*$ である．
固定した $x$ に対して $y\mapsto b(x,y)$ は
有界線形汎関数だから
\[
 \sup_{\|y\|\le1}|T_y(x)|<\infty
\]
である．定理\ref{thm:ubp}を族
$\{T_y:\|y\|\le1\}$ に適用して
$\sup_{\|y\|\le1}\|T_y\|=C<\infty$ を得る．
よって $\|y\|\le1$ なら
$|b(x,y)|\le C\|x\|$ である．
$y\ne0$ について $y/\|y\|$ を代入し，
$y=0$ では双線形性を用いれば結論が従う．
最後に
\[
 b(x+h,y+k)-b(x,y)
 =b(h,y)+b(x,k)+b(h,k)
\]
を上の評価に代入すれば連続性を得る．
\end{proof}

\begin{proposition}[作用素族の同程度連続性]\label{prop:equicont}
ノルム空間 $X,Y$ と
$(T_\lambda)_{\lambda\in\Lambda}\subset\mathcal B(X,Y)$ に対し，
次は同値である．
\begin{enumerate}
\item $\sup_{\lambda\in\Lambda}\|T_\lambda\|<\infty$．
\item 任意の $\varepsilon>0$ に対しある $\delta>0$ が存在して，
      全ての $\lambda$ と $\|x\|<\delta$ に対し
      $\|T_\lambda x\|<\varepsilon$ が成立する．
\end{enumerate}
特に $X$ がBanach空間なら，この二条件は
各点有界性とも同値である．
\end{proposition}
\begin{proof}
$M=\sup_\lambda\|T_\lambda\|<\infty$ とすれば，
$\delta=\varepsilon/(M+1)$ と取ることにより
$\|T_\lambda x\|\le M\|x\|<\varepsilon$ を得る．
逆に第2条件を $\varepsilon=1$ で適用し，
対応する $\delta>0$ を取る．
$\|x\|\le1$ なら
$\|\delta x/2\|<\delta$ だから，
全ての $\lambda$ に対して
\[
 \|T_\lambda x\|
 =\frac{2}{\delta}
   \left\|T_\lambda\!\left(\frac{\delta x}{2}\right)\right\|
 <\frac{2}{\delta}
\]
である．従って $\sup_\lambda\|T_\lambda\|\le2/\delta$ となる．
最後の主張は一様有界性原理と，
一様な作用素ノルムの上界が各点有界性を導くことから従う．
\end{proof}

\section{特異性の凝縮}

一様有界性原理の対偶は，
作用素ノルムが一様有界でない作用素族に対し，
ある点での像が非有界となることを意味する．
実際には，そのような点は位相的に大きな集合をなす．

\begin{theorem}[共鳴の原理]\label{thm:condensation}
$X$ をBanach空間，$Y$ をノルム空間，
$(T_\lambda)\subset\mathcal B(X,Y)$ とする．
\[
 \sup_{\lambda\in\Lambda}\|T_\lambda\|=\infty
\]
であるなら，
\[
 E=\left\{x\in X:
  \sup_{\lambda\in\Lambda}\|T_\lambda x\|=\infty\right\}
\]
は稠密な $G_\delta$ 集合である．
\end{theorem}
\begin{proof}
各 $n\ge1$ に対して
\[
 F_n=\{x\in X:\sup_\lambda\|T_\lambda x\|\le n\}
\]
と置く．$F_n$ が閉であることは
定理\ref{thm:ubp}の証明のとおりである．
もしある $F_n$ が内部を持つなら，
同じ証明の球上の評価から
$\sup_\lambda\|T_\lambda\|<\infty$ が従い，
仮定に矛盾する．
従って各 $F_n$ は至る所疎な閉集合である．
そこで
\[
 E=X\setminus\bigcup_{n=1}^\infty F_n
   =\bigcap_{n=1}^\infty(X\setminus F_n)
\]
は稠密開集合の可算共通部分である．
Baireのカテゴリー定理によって $E$ は稠密であり，
定義から $G_\delta$ 集合でもある．
\end{proof}

\begin{corollary}
ノルムが非有界な作用素列 $(T_n)$ について，
$\sup_n\|T_nx\|=\infty$ を満たす $x$ は，
任意の空でない開集合の中に存在する．
\end{corollary}
\begin{proof}
定理\ref{thm:condensation}による集合 $E$ の稠密性を用いる．
\end{proof}

\section{完備性の必要性と反例}

\begin{example}[稠密集合上の各点有界性は不十分]
$X=\ell^2(\mathbb K)$ とし，
$T_n:X\to\mathbb K$ を $T_nx=nx_n$ と定める．
このとき各 $T_n$ は連続であり，$\|T_n\|=n$ である．
しかし，稠密な線形部分空間 $c_{00}\subset\ell^2$ 上では
族 $(T_n)$ は各点有界である．
\end{example}
\begin{proof}
$|T_nx|=n|x_n|\le n\|x\|_2$ より $\|T_n\|\le n$，
標準基底 $e_n$ について $T_ne_n=n$ より
$\|T_n\|=n$ となる．
任意の $x\in c_{00}$ は有限台を持つため
$\sup_n|nx_n|<\infty$ である．
$c_{00}$ の稠密性は，任意の $x=(x_j)\in\ell^2$ を
その最初の $N$ 成分だけで切った元 $x^{(N)}$ により
\[
 \|x-x^{(N)}\|_2^2=\sum_{j>N}|x_j|^2\to0
\]
となることから従う．
一方 $x_j=j^{-3/4}$ と置くと
$\sum_j|x_j|^2=\sum_jj^{-3/2}<\infty$ であるのに，
$T_nx=n^{1/4}\to\infty$ となる．
従って完備な空間の稠密部分集合上で各点有界でも，
空間全体で各点有界とは限らない．
\end{proof}

\begin{example}[非完備空間での失敗]\label{ex:c00}
$c_{00}$ を有限個を除き全ての成分が $0$ である
数列全体とし，$\ell^2$ ノルムを入れる．
$T_n:c_{00}\to\mathbb K$ を
\[
 T_nx=n x_n
\]
と定めると，各 $T_n$ は有界線形汎関数で，
族 $(T_n)$ は各点有界だが，
$\sup_n\|T_n\|=\infty$ である．
\end{example}
\begin{proof}
各 $n$ について
$|T_nx|=n|x_n|\le n\|x\|_2$ なので
$\|T_n\|\le n$ である．標準基底 $e_n\in c_{00}$ により
$T_ne_n=n$ だから $\|T_n\|=n$ である．
一方，任意の $x\in c_{00}$ では
$x_n=0$ は十分大きい $n$ について成立するので
$\sup_n|nx_n|<\infty$ となる．

$c_{00}$ が完備でないことを直接確認する．
$x=(2^{-n})_{n\ge1}\in\ell^2\setminus c_{00}$ とし，
その最初の $m$ 成分以外を $0$ にした数列を
$x^{(m)}\in c_{00}$ とすると
\[
 \|x^{(m)}-x^{(k)}\|_2^2
 \le\sum_{n>\min\{m,k\}}4^{-n}\longrightarrow0
\]
である．従って $(x^{(m)})$ は $c_{00}$ のCauchy列だが，
$c_{00}$ 内では収束しない．
実際，収束するなら各座標の極限は $2^{-n}$ でなければならず，
その極限は $c_{00}$ に属さない．
\end{proof}

\begin{example}[非完備空間上の不連続な各点極限]
同じ $c_{00}$ について
\[
 S_nx=\sum_{j=1}^n jx_j
\]
と置く．各 $S_n$ は有界線形汎関数で，
任意の $x\in c_{00}$ について $S_nx$ は
\[
 Sx=\sum_{j=1}^\infty jx_j
\]
に収束するが，$S$ は有界でない．
\end{example}
\begin{proof}
有限和にCauchy--Schwarzの不等式を適用すると
\[
 |S_nx|\le\left(\sum_{j=1}^n j^2\right)^{1/2}\|x\|_2
\]
となるので $S_n$ は有界である．
$c_{00}$ の各元は有限台なので，その台を含む $n$ 以降では
$S_nx=Sx$ となる．
しかし $S(e_j)=j$ かつ $\|e_j\|_2=1$ だから，
$S$ が有界だと仮定すると
$j\le\|S\|$ がすべての $j$ で成立して矛盾する．
\end{proof}


\section{到達課題の完全解答}
\noindent\textbf{Banach--Steinhausの完備性を検証する．}
\begin{proposition}
$c_{00}$ を有限個の成分しか非零でない実数列の空間とし，
$\|x\|_\infty=\sup_k|x_k|$ とする．
線形汎関数 $T_n(x)=nx_n$ は各 $x\in c_{00}$ について
$\sup_n|T_n(x)|<\infty$ を満たすが，
$\sup_n\|T_n\|=\infty$ である．
\end{proposition}
\begin{proof}
$x$ は有限個の座標以外では零だから，
十分大きい $n$ で $T_n(x)=0$ となり，
値の集合は有界である．
一方 $|T_n(x)|\le n\|x\|_\infty$ であり，
第 $n$ 単位ベクトル $e_n$ に対し
$T_n(e_n)=n$，$\|e_n\|_\infty=1$ だから
$\|T_n\|=n$ となる．
$c_{00}$ は $\ell^\infty$ の閉部分空間ではなく，
例えば $(1/n)_{n\ge1}$ は有限支持列によって
一様近似できるが $c_{00}$ には属さない．
したがって $c_{00}$ はBanach空間でなく，
一様有界性原理の仮定を満たさない．
\end{proof}
\noindent\textbf{到達確認．}
$\ell^\infty$ 上の $T_n(x)=x_n$ の作用素ノルムは幾つか．
\par\noindent\textbf{解答．}
$|x_n|\le\|x\|_\infty$ で，上界1．
$e_n$ により等号を達成するので1である．
\end{document}
